% =============================================================================
%  Abstract_Template.tex — the 3-sentence abstract exercise, standalone
%  Mastering the Scientific Narrative · Deep Learning Indaba 2026
%
%  A dependency-free extract of the abstract exercise from
%  indaba-narrative-template.zip (main.tex), for anyone who wants just the
%  writing scaffold without the full paper skeleton or the indaba_narrative
%  package. Compiles standalone with a plain \documentclass{article}.
%
%  Rules:
%   * Sentence 1 (Context + Gap) must contain a place and a number.
%   * Sentence 2 (Method) must name the thing you built and what is new.
%   * Sentence 3 (Results + Impact) must contain a metric, a value, and a
%     NAMED baseline.
%   * No sentence may contain "very", "novel" or "promising" without a
%     number next to it.
% =============================================================================
\documentclass[11pt]{article}
\usepackage[a4paper,margin=2.4cm]{geometry}
\usepackage[utf8]{inputenc}
\usepackage[T1]{fontenc}

\title{\bfseries [Working title: Method]: [what it does] for [problem]}
\author{Your Name \and Co-author Name}
\date{\today}

\begin{document}
\maketitle

\begin{abstract}
% --- S1 · CONTEXT + GAP -----------------------------------------------------
% Where, who, how big is the stake — and what blocks the world today?
% e.g. "In West Africa, cassava mosaic and brown-streak disease destroy up
% to 40% of yields for smallholder farmers, yet existing vision-based
% diagnostic models assume cloud connectivity and high-end hardware that
% rural farms do not have."
\textbf{[S1 — Context + Gap:]} \dots

% --- S2 · METHOD -------------------------------------------------------------
% Your one-line answer: what did you build, and what is new about it?
% e.g. "We introduce CropGuard, a quantized and pruned convolutional
% detector that runs entirely offline on low-end Android devices, trained
% on a new dataset of 12,406 field images collected across 42 farms in
% Guinea and Nigeria and released under FAIR principles."
\textbf{[S2 — Method:]} \dots

% --- S3 · RESULTS + IMPACT ---------------------------------------------------
% Named metrics, exact numbers, named baselines — then the door it opens.
% e.g. "CropGuard reaches 78.3% top-1 accuracy while cutting inference
% latency by 40% relative to MobileNetV3, enabling farmers to diagnose
% disease in the field with no internet connection."
\textbf{[S3 — Results + Impact:]} \dots
\end{abstract}

% Self-check before you move on:
%  [ ] A stranger could FUND this project after these 3 sentences.
%  [ ] Sentence 3 contains at least one NUMBER.
%  [ ] The gap (what blocks the world today) is explicit.
%  [ ] The region/context is named — an anchor, not a footnote.
%  [ ] Zero empty adjectives ("good", "promising", "novel") survive.

\end{document}
